\section{Bayesian Framework}

\subsection{Prior and Posterior}

Let $\pi^*$ be the true percentile of the enacted plan in the distribution
of all valid redistricting plans. We assume a uniform prior:
$\pi^* \sim \mathrm{Uniform}(0, 1)$.

The ensemble provides an estimate $\hat{p}$ of $\pi^*$, based on a sample
of $N$ plans with effective sample size $n^* = \mathrm{ESS}$.
Under the Beta-Binomial model, the likelihood of observing $k = \lfloor \hat{p} \cdot N \rfloor$
plans more compact than the enacted plan (out of $N$ total) is:
\[
  P(k \mid \pi^*) = \binom{N}{k} (\pi^*)^k (1-\pi^*)^{N-k}.
\]

With a uniform prior ($\mathrm{Beta}(1,1)$), the posterior is:
\[
  \pi^* \mid k, N \;\sim\; \mathrm{Beta}(k+1, N-k+1).
\]

In practice we use $n^*$ in place of $N$ to account for autocorrelation:
\[
  \pi^* \mid \hat{p}, n^* \;\sim\; \mathrm{Beta}(\hat{p} \cdot n^* + 1,\;
  (1-\hat{p})\cdot n^* + 1).
\]

\subsection{The Bayesian Detection Score}

\begin{definition}[Bayesian Detection Score]
\label{def:bds}
The BDS at threshold $\tau$ is:
\[
  \mathrm{BDS}(\tau) = P(\pi^* < \tau \mid \hat{p}, n^*)
  = I_\tau(\hat{p}\cdot n^* + 1,\; (1-\hat{p})\cdot n^* + 1),
\]
where $I_\tau(a, b)$ is the regularised incomplete Beta function (the CDF of
$\mathrm{Beta}(a, b)$ evaluated at $\tau$).
We use $\tau = 0.05$ (the plan is in the most compact 5\%) as the default
threshold.
\end{definition}

BDS$(\tau) > 0.95$ means the posterior probability exceeds 95\% that the plan
is genuinely in the most compact $\tau$-fraction. This is analogous to $p < 0.05$
in frequentist testing, but incorporates the ensemble's finite size via $n^*$.

\subsection{The Geography Defense and Bayesian Reasoning}

The geography defense argues that the enacted plan's compactness is explained
by geographic sorting, not partisan intent. In the Bayesian framework, the
geography defense is absorbed into the \emph{prior}: the ensemble is constructed
by running ReCom on the actual census-tract graph, which already incorporates
all geographic constraints. The posterior therefore measures only the residual
above and beyond what geography predicts.

A BDS near 1.0 means the plan is extreme even within the geography-constrained
space of valid plans — the geography defense does not explain the plan's position.
A BDS near 0.0 means the plan is consistent with a geographically-constrained
neutral draw.

\subsection{Comparison to Frequentist Percentile}

The frequentist percentile $\hat{p}$ is a point estimate with no uncertainty
quantification. The BDS uses the posterior distribution to incorporate:
\begin{enumerate}
  \item The uncertainty from finite $N$ (small ensembles have wide posteriors)
  \item The ESS correction (autocorrelated chains give less information than
    their nominal $N$ suggests)
\end{enumerate}
For large, well-mixed ensembles ($n^* \gg 1/\hat{p}$), the BDS converges to
an indicator function ($\approx 1$ if $\hat{p} < \tau$, $\approx 0$ if
$\hat{p} > \tau$). For small ensembles (typical in current practice), the BDS
provides a meaningful probability rather than a threshold answer.
